\documentclass[letterpaper]{article} % DO NOT CHANGE THIS
\usepackage{aaai2027}  % DO NOT CHANGE THIS
% The serif, sans-serif, and monospaced fonts are loaded automatically by
% aaai2027.sty (newtxtext, helvet, courier). DO NOT add \usepackage{times},
% \usepackage{helvet}, \usepackage{courier}, or any other font package.
\usepackage[hyphens]{url}  % DO NOT CHANGE THIS
\usepackage{graphicx} % DO NOT CHANGE THIS
\urlstyle{rm} % DO NOT CHANGE THIS
\def\UrlFont{\rm}  % DO NOT CHANGE THIS
\usepackage{natbib}  % DO NOT CHANGE THIS AND DO NOT ADD ANY OPTIONS TO IT
\usepackage{caption} % DO NOT CHANGE THIS AND DO NOT ADD ANY OPTIONS TO IT
\frenchspacing  % DO NOT CHANGE THIS
%
% These are recommended to typeset algorithms but not required. See the subsubsection on algorithms. Remove them if you don't have algorithms in your paper.
\usepackage{algorithm}
\usepackage{algorithmic}

%
% These are recommended to typeset listings but not required. See the subsubsection on listing. Remove this block if you don't have listings in your paper.
\usepackage{newfloat}
\usepackage{listings}
\DeclareCaptionStyle{ruled}{labelfont=normalfont,labelsep=colon,strut=off} % DO NOT CHANGE THIS
\lstset{%
	basicstyle={\footnotesize\ttfamily},% footnotesize acceptable for monospace
	numbers=left,numberstyle=\footnotesize,xleftmargin=2em,% show line numbers, remove this entire line if you don't want the numbers.
	aboveskip=0pt,belowskip=0pt,%
	showstringspaces=false,tabsize=2,breaklines=true}
\floatstyle{ruled}
\newfloat{listing}{tb}{lst}{}
\floatname{listing}{Listing}

%
% Recommended for better-looking tables
\usepackage{booktabs}

%
% Keep the \pdfinfo as shown here. There's no need
% for you to add the /Title and /Author tags.
\pdfinfo{
/TemplateVersion (2027.1)
}

% DISALLOWED PACKAGES
% \usepackage{authblk} -- This package is specifically forbidden
% \usepackage{balance} -- This package is specifically forbidden
% \usepackage{CJK} -- This package is specifically forbidden
% \usepackage{float} -- This package is specifically forbidden
% \usepackage{flushend} -- This package is specifically forbidden
% \usepackage{fullpage} -- This package is specifically forbidden
% \usepackage{geometry} -- This package is specifically forbidden
% \usepackage{hyperref} -- This package is specifically forbidden
% \usepackage{navigator} -- This package is specifically forbidden
% (or any other package that embeds links such as navigator or hyperref)
% \usepackage{indentfirst} -- This package is specifically forbidden
% \usepackage{layout} -- This package is specifically forbidden
% \usepackage{multicol} -- This package is specifically forbidden
% \usepackage{nameref} -- This package is specifically forbidden
% \usepackage{savetrees} -- This package is specifically forbidden
% \usepackage{setspace} -- This package is specifically forbidden
% \usepackage{stfloats} -- This package is specifically forbidden
% \usepackage{tabu} -- This package is specifically forbidden
% \usepackage{titlesec} -- This package is specifically forbidden
% \usepackage{tocbibind} -- This package is specifically forbidden
% \usepackage{ulem} -- This package is specifically forbidden
% \usepackage{wrapfig} -- This package is specifically forbidden
% DISALLOWED COMMANDS
% \nocopyright -- Your paper will not be published if you use this command
% \addtolength -- This command may not be used
% \balance -- This command may not be used
% \baselinestretch -- Your paper will not be published if you use this command
% \clearpage -- No page breaks of any kind may be used for the final version of your paper
% \columnsep -- This command may not be used
% \newpage -- No page breaks of any kind may be used for the final version of your paper
% \pagebreak -- No page breaks of any kind may be used for the final version of your paper
% \pagestyle -- This command may not be used
% \tiny -- This is not an acceptable font size.
% \vspace{- -- No negative value may be used in proximity of a caption, figure, table, section, subsection, subsubsection, or reference
% \vskip{- -- No negative value may be used to alter spacing above or below a caption, figure, table, section, subsection, subsubsection, or reference

\setcounter{secnumdepth}{0} %May be changed to 1 or 2 if section numbers are desired.

% The file aaai2027.sty is the style file for AAAI Press
% proceedings, working notes, and technical reports.
%

% Title
% Single-blind submission: author names appear on the paper.
\title{PITBULL: Temporal Safety Skill for Coding Agents}
\author{
    Pavel Marian\textsuperscript{\rm 1},
    Stanislav Chumakov\textsuperscript{\rm 1},
    Danil Ezhov\textsuperscript{\rm 1},
    Alexander Boukhanovsky\textsuperscript{\rm 1}
}
\affiliations{
    \textsuperscript{\rm 1}AI Institute, ITMO University, Saint Petersburg, Russian Federation\\
    svchumakov@itmo.ru, doezhov@itmo.ru, boukhanovsky@mail.ifmo.ru
}

\begin{document}

\maketitle

\begin{abstract}
We present \textsc{Pitbull}, an executable temporal-safety skill for
agent-generated feature programs. The skill detects point-in-time violations
by comparing pipeline outputs on full and temporally truncated databases,
adds canary perturbations for late-arriving fields, and validates every
agent repair on held-out prediction times. Across three runs over 62 leaking
programs, 97.8\% of program--run pairs were repaired within five iterations.
Packaged as a skill that an independent coding agent selects on its own, the
same loop repairs 76 of 76 programs, matching direct invocation of the
procedure.
\end{abstract}

\begin{links}
    \link{Demo}{https://pavelmarian.github.io/pitbull/demo/temporal-safety/}
    \link{Code}{https://github.com/pavelmarian/pitbull}
\end{links}

\section{Introduction and Related Work}

Successful execution of a relational feature pipeline by a coding agent
does not guarantee its temporal correctness. A pipeline that filters on a
single reference timestamp can still read information that did not exist at
the prediction time, because availability is spread across related tables
and columns. Such violations yield plausible features and inflated offline
results that do not reproduce at serving time. The property that every
feature computed at time $t$ depends only on facts available at time $t$ is
called point-in-time correctness.

Temporal leakage is a recurring source of irreproducible machine-learning
results~\cite{10.1145/2382577.2382579,kapoor2022leakagereproducibilitycrisismlbased},
including in temporally split process-mining
benchmarks~\cite{DBLP:journals/corr/abs-2107-01905}. The relational setting
makes it harder: features are computed across several
tables~\cite{robinson2024relbenchbenchmarkdeeplearning}, and the availability
of a value may be governed by another row or by a column-specific timestamp.
Static notebook analysis and pipeline instrumentation target preprocessing
and train--test
contamination~\cite{yang2022dataleakagenotebooksstatic,10.1145/3448016.3452759}.
LLM-based code inspection is probabilistic~\cite{alturayeif2025data} and
insufficient as the only safeguard.

The \textsc{Pitfall} system introduced black-box differential execution for
detecting and localising point-in-time
violations~\cite{chumakov2025pitfall}. \textsc{Pitbull} extends it from
diagnosis to a closed agentic loop of detection, repair, and independent
revalidation. The verification is one-sided, in line with the view of
look-ahead freedom as temporal non-interference: an observed divergence is a
finite witness of a violation, and the absence of divergence at finitely
many probes is no
proof~\cite{fonseca2026lookaheadfreedomtemporalnoninterferenceverifiable}.

\begin{figure*}[t]
    \centering
    \includegraphics[width=0.8\linewidth]{Frame 8 (2).png}
    \caption{\textsc{Pitbull} architecture: the agent submits a feature
    program; the skill runs it on the full, the truncated, and the
    canary-perturbed database, reports diverging features, and gates each
    repair on development and held-out prediction times.}
    \label{fig:system_architecture}
\end{figure*}

\section{System Overview}

\textsc{Pitbull} is a skill that a coding agent invokes over a candidate
feature program $\varphi(D, e, t)$: a database $D$, entities $e$, and a
prediction time $t$. The agent writes and revises the program; the skill
only executes it and returns evidence (Figure~\ref{fig:system_architecture}).

\textbf{Differential witness.}
For each prediction time the skill runs the program twice, on the full
database and on the database truncated at $t$:
\begin{equation}
  y_{\mathrm{full}} = \varphi(D,e,t), \qquad
  y_{\mathrm{past}} = \varphi(D|_t,e,t).
\end{equation}
Any difference beyond a numerical tolerance is a \texttt{LEAK} witness, and
the skill names the affected output columns. Equality only says that no
violation was observed at the times tested.

\textbf{Canary.}
Some channels are hidden by aggregation or exercised weakly by real future
values. The skill therefore builds a counterfactual future: numbers are
shifted, categories replaced by a reserved value, dates moved beyond the
data horizon, including late-filled fields such as delivery dates. An
output that reacts to this perturbation is reported as a canary hit. Canary
is less specific than the witness (it fired on 4 of 30 clean control
programs), so the report lists canary columns separately for a human to
judge.

\textbf{Repair loop and held-out gate.}
On a \texttt{LEAK}, the agent rewrites the program under a fixed
natural-language instruction that names the class of error (features must
use only information that exists at $t$; reviews and delivery dates arrive
after the purchase) and reruns the check. A candidate is \texttt{CLEAN} only
if it executes and shows neither witness nor canary divergence on all
development and held-out prediction times; missing code, runtime errors and
timeouts count as failures. The loop stops at the first \texttt{CLEAN} or
after five repairs. The full history of candidates and their verdicts is
kept in a state file next to the code. A \texttt{CLEAN} verdict is scoped to
the declared availability map and the tested times.

\section{Demonstration}

The demonstration runs in the browser at the linked page. Visitors pick an
Olist profile and a candidate feature program (the bundled sample or their
own file), set a development prediction time, and press \emph{Run temporal
check}. The page shows the three data views side by side, the diverging
product features, the availability channel that caused the divergence (for
the sample, reviews joined without a \texttt{review\_creation\_date} cut),
and the acceptance gate on development and held-out times. A repair by the
coding agent is then applied, re-checked, and offered for download as a
patched source, a diff, and a JSON evidence report. Four study panels let
visitors browse the evidence behind Table~\ref{tab:results}: CLEAN@$k$ across
the corpus, a before/after view of one repair, the three repeated runs, and
the number of prediction times needed to expose a leak. A live run takes
under a minute per iteration. The same skill loads into a coding agent
unchanged (Table~\ref{tab:results}, last row).

\section{Experimental Snapshot}

Table~\ref{tab:results} summarises repair outcomes on leaking relational
feature programs derived from Olist tasks, with one repair model. In three
independent direct runs over 62 programs, 182 of 186 program--run pairs
became clean within five iterations. A density experiment shows that two
prediction times give 97.6\% conservative witness detection; the witness had
no false positives on 30 clean control programs, the canary fired on four.

\begin{table}[t]
\centering\footnotesize
\setlength{\tabcolsep}{2pt}
\begin{tabular}{@{}lrrrrr@{}}
\toprule
Invocation & $n$ & CLEAN@1 & CLEAN@5 & Calls & Cost \\
\midrule
Direct, 1 run        & 76          & 85.5\% & 98.7\% & 91  & \$0.12 \\
Direct, 3 runs       & 62$\times$3 & 84.4\% & 97.8\% & 255 & \$0.32 \\
Agent-selected skill & 76          & 98.7\% & 100\%  & 578 & \$0.48 \\
\bottomrule
\end{tabular}
\caption{Repair outcomes with the same repair model. The skill row uses the
same 76 programs as the first row; exact McNemar over paired per-program
outcomes: 0 vs.\ 1 discordant, $p = 1.0$.}
\label{tab:results}
\end{table}

We then packaged the loop as a reusable skill and ran it under an
independent agentic harness that selects and invokes the skill on its own;
the skill name does not appear in the task prompt. On the same 76 programs
the skill matches direct invocation: CLEAN@5 reaches 100\%, and the one
discordant program is the single failure of direct invocation, repaired at
the first iteration under the skill. In skill mode the agent also sees the
checker output with column names, which explains part of the gain at
CLEAN@1. The price is packaging overhead: about $4\times$ the inference cost
and $6\times$ the model calls per program, driven by the harness system
prompt on every turn. The evidence is strongest for review-based leakage,
which dominates the corpus.

\section{Conclusion}

\textsc{Pitbull} gives a coding agent an executable way to detect, localise,
repair, and revalidate point-in-time violations in relational feature
pipelines, with the verdict made by the checker and the repair by the agent.
Future work: independently annotated per-column availability maps, a corpus
beyond review-based leakage, and automatic inference of availability maps
under the same evidentiary standard.

\bibliography{aaai2027}

\end{document}
